% =============================================================================
%  2 (cont.).  MODEL SECTIONS MISSING FROM 02-model.tex
%
%  Written against code/global/ as of 2026-08-31 (branch bocola-rewrite).
%  Notation, macros (\D, \F, \E) and style follow 02-model.tex; nothing in that
%  file is edited here.  Each block below carries an INSERT marker saying where
%  it belongs.
%
%  NEW SYMBOLS INTRODUCED (none collide with 02-model.tex or 03-calibration.tex):
%    \mathsf{M}_X   country mass            \varkappa   size ratio M_F/M_D = 8
%    V_t            cross-border deposit position (D-good units)
%    W_{X,t}        household gross carried deposit claim
%    \kappa_{nfa}   SGU debt-elastic premium on the external position
%    \psi_X         cross-border sovereign portfolio adjustment cost
%    \pi^{\ell}     per-period probability the CB facility is offered
%    \ell_X         size of the facility        \mathsf{cb}_t  facility indicator
%    \rho_z         TFP persistence         \mathcal{S}_t  the aggregate state
%    \Gamma_X       cross-sectional distribution over (a,e)
%    \mathcal{T}_X  the aggregation constant in the representative-agent budget
%
%  BIB KEYS REQUIRED BEYOND THOSE ALREADY USED:
%    schmittgrohe2003closing, carroll2006method, young2010solving,
%    krusell1998income, smolyak1963, kruegerkubler2004, judd2014smolyak,
%    gertlerkiyotaki2010 (already in 03-calibration.tex)
%
%  ONE SUBSTANTIVE CORRECTION TO 02-model.tex IS PROPOSED, NOT APPLIED:
%    eq. (bohn) states the fiscal rule as a constant ELASTICITY.  The solved
%    model uses a rule LINEAR in the level of the surviving stock whose
%    coefficient is inverted from a target debt root, with a REGIME-DEPENDENT
%    anchor.  See \cref{sec:fiscal-rule-corrected} for the statement, the
%    derivation and the reason the elasticity form fails here.
% =============================================================================


% =============================================================================
% INSERT: immediately after the opening paragraph of \section{The Model},
%         before \subsection{Households}.
% =============================================================================
\subsection{Country size and the per-capita convention}
\label{sec:size}

The two countries are not of equal size. Country $X$ is populated by a continuum
of households of mass $\mathsf{M}_X$, and we write
\begin{equation}
\varkappa\;\equiv\;\frac{\mathsf{M}_\F}{\mathsf{M}_\D}
\label{eq:size-ratio}
\end{equation}
for the size of the core relative to the periphery. \emph{Every} variable in this
paper --- output, hours, capital, net worth, deposits, bond holdings --- is
expressed \emph{per capita of its own country}, so that $\varkappa$ appears in
exactly one place: wherever a $\D$ quantity and an $\F$ quantity are added
together. There are four such places, and they are the goods market
\eqref{eq:mkt-goods}, the union deposit market \eqref{eq:mkt-dep}, the two
sovereign markets \eqref{eq:mkt-bD}--\eqref{eq:mkt-bF}, and the union wealth
identity \eqref{eq:union-wealth}. Setting $\varkappa=1$ reproduces the symmetric
two-country model exactly.

Asymmetric size is not a cosmetic realism device; it is what keeps the sovereign
shock from being arbitraged away by the union deposit market. In a symmetric
union the periphery is half of the union, so $\D$'s own sovereign shock moves the
union-wide real deposit rate --- by $45$ basis points per year at the headline
experiment --- and that fall in $r_\D$ cancels roughly four-fifths of the rise in
the credit spread $\lambda_\D\mu_\D/\E_t[\Omega_\D]$ before it ever reaches a
firm's wage bill through \eqref{eq:rwc}. \citet{bocola2016pass}'s own
open-economy extension has no such feedback because his intermediaries face a
\emph{world} interest rate, exogenous to the country in crisis. Making $\F$ large
is the general-equilibrium counterpart of that assumption: it is the smallest
departure from a genuine union that restores the exogeneity of the union rate to
$\D$'s own shock, and it nests his environment as $\varkappa\to\infty$.

\paragraph{Home bias must scale with size.} With per-capita quantities and
unequal masses, a single home-bias parameter is inconsistent with balanced trade.
Read $\varpi$ in \eqref{eq:ces-D} as $\varpi_\D$ and in \eqref{eq:ces-F} as
$\varpi_\F$. At $p=1$ and equal per-capita consumption, $\D$'s aggregate imports
are $\mathsf{M}_\D(1-\varpi_\D)C$ and $\F$'s are $\mathsf{M}_\F(1-\varpi_\F)C$,
so trade balances if and only if
\begin{equation}
1-\varpi_\F\;=\;\frac{1-\varpi_\D}{\varkappa},
\label{eq:home-bias-size}
\end{equation}
which is imposed as a restriction rather than calibrated: $\varpi_\F$ is
\emph{derived} from $\varpi_\D$ and $\varkappa$. The small country is the open
one. At $\varpi_\D=0.85$ and $\varkappa=8$, $\D$ imports $15\%$ of its basket and
$\F$ imports $1.875\%$ of its.

\paragraph{Bilateral flows.} Let $IM_{\D,t}$ be $\D$'s per-capita import of
$\F$-goods, in $\F$-good units, and $IM_{\F,t}$ be $\F$'s per-capita import of
$\D$-goods, in $\D$-good units, both given by \eqref{eq:ces-D}--\eqref{eq:ces-F}.
Each country's exports are consumed by the \emph{other} country's population, so
net exports in own-good units per capita are
\begin{equation}
NX_{\D,t}=\varkappa\,IM_{\F,t}-p_t\,IM_{\D,t},
\qquad
NX_{\F,t}=\frac{IM_{\D,t}}{\varkappa}-\frac{IM_{\F,t}}{p_t}.
\label{eq:nx}
\end{equation}

\paragraph{Sovereign holdings.} A bond position is a claim on one issuer held by
two different populations, so its units must be fixed once. Home holdings are
carried in the holder's own per-capita units; \emph{both cross-border legs are
carried per $\D$ capita}. Thus $b^{\D}_{\F,t+1}$ denotes the $\F$ banking
sector's aggregate holding of the $\D$ sovereign divided by $\mathsf{M}_\D$, so
that the $\F$ bank's own per-capita book is $b^{\D}_{\F,t+1}/\varkappa$, and
market clearing in the $\D$ sovereign is the unweighted sum
\eqref{eq:mkt-bD}. Symmetrically $b^{\F}_{\D,t+1}$ is $\D$'s own per-capita
holding of the $\F$ sovereign, which enters the $\F$ market divided by
$\varkappa$. Under \eqref{eq:home-bias-size} and this convention the per-capita
steady state is \emph{identical} to the symmetric one: $\bar p=1$, the same
leverage, the same net worth, the same deposit supply in both countries. Size
changes the model's general equilibrium, not its steady state.


% =============================================================================
% INSERT: as the closing part of \subsection{Households}.
% =============================================================================
\subsection{Idiosyncratic risk, aggregation, and what the household block does}
\label{sec:hh-aggregation}

\paragraph{The income process.} Idiosyncratic labour productivity $e_{it}$ in
\eqref{eq:hh-budget} follows an AR(1) in logs,
$\log e_{i,t+1}=\rho_e\log e_{it}+\sigma_e\varepsilon^{e}_{i,t+1}$, discretised by
the \citet{rouwenhorst1995} method into an $n_e$-state Markov chain
$(\mathsf{e},\Pi_X)$ whose grid is normalised so that
$\sum_e\pi^{\ast}(e)\mathsf{e}(e)=1$ under the chain's stationary distribution
$\pi^{\ast}$. Idiosyncratic risk is uninsurable, the only asset is the deposit
$a_{it}$, and the borrowing limit \eqref{eq:hh-borrowing} binds for a positive
mass of households. Let $\Gamma_{X,t}(a,e)$ denote the cross-sectional
distribution and $a'_X(a,e)$ the savings policy implied by
\eqref{eq:hh-objective}--\eqref{eq:hh-borrowing}. Aggregate saving and
consumption are
\begin{equation}
A_{X,t}=\int a'_X(a,e)\,\mathrm{d}\Gamma_{X,t},
\qquad
C_{X,t}=\int c_X(a,e)\,\mathrm{d}\Gamma_{X,t},
\qquad
\Gamma_{X,t+1}=\mathcal{H}_X\bigl[\Gamma_{X,t}\bigr],
\label{eq:hh-aggregation}
\end{equation}
with $\mathcal{H}_X$ the forward operator induced by $a'_X$ and $\Pi_X$. Policies
are computed by the endogenous grid method \citep{carroll2006method} on the GHH
composite, and $\mathcal{H}_X$ by the lottery method of
\citet{young2010solving}.

\paragraph{What the household block is for, and what closes the aggregate model.}
The incomplete-markets block plays two roles here, and it is worth being explicit
that neither of them is the propagation of aggregate shocks. First, it delivers
the economy's \emph{deposit-supply schedule}: the discount factor $\beta_X$ is
not calibrated to a wealth target but solved so that
\eqref{eq:hh-aggregation} clears the intermediaries' funding need at the
calibrated deposit rate, $A_X(\beta_X;\bar r)=\overline{\mathrm{dep}}_X$, which is
what makes the steady-state supply of bank liabilities an equilibrium object
rather than an assumption. Second, it supports the distributional accounting of
the policy experiment: households are run along the solved aggregate paths to
compute consumption-equivalent welfare by income quintile.

The \emph{aggregate} dynamics, however, close with the representative-agent
counterpart of \eqref{eq:hh-aggregation}. Aggregate consumption is the residual of
the aggregate budget and the deposit rate is pinned by an aggregate Euler
equation on the GHH composite $x_{X,t}\equiv C_{X,t}-v(N_{X,t})$,
\begin{align}
C_{X,t}&=\frac{W_{X,t}}{P^{c}_{X,t}}
 +\frac{w_{X,t}N_{X,t}+\mathrm{Div}_{X,t}-T^{\tau}_{X,t}}{P^{c}_{X,t}}
 +\mathcal{T}_X-A_{X,t},
\label{eq:agg-budget}\\[3pt]
x_{X,t}^{-\sigma_X}
&=\beta^{\rm eff}_X\,
  \E_t\Bigl[\bigl(1+r^{c}_{X,t+1}\bigr)\,x_{X,t+1}^{-\sigma_X}\Bigr],
\qquad
\beta^{\rm eff}_X\equiv\frac{1}{1+\bar r_X},
\label{eq:agg-euler}
\end{align}
where $W_{X,t}$ is the household's gross carried deposit claim
(\cref{sec:union-deposits}), $1+r^{c}_{X,t+1}=(1+r_{X,t})P^{c}_{X,t}/P^{c}_{X,t+1}$
is the predetermined real deposit return, and $\mathcal{T}_X$ is a constant chosen
once so that \eqref{eq:agg-budget} returns exactly the incomplete-markets
aggregates $(\bar C_X,\bar A_X)$ at the steady state. Two things are being said
here. The choice $\beta^{\rm eff}_X=1/(1+\bar r_X)$ makes \eqref{eq:agg-euler}
reproduce the heterogeneous-agent aggregate at the steady state exactly ---
incomplete markets require $\beta_X(1+\bar r_X)<1$, and $\mathcal{T}_X$ together
with $\beta^{\rm eff}_X$ absorbs the precautionary wedge as a level shift. And
the constant $\mathcal{T}_X$ is an aggregation constant, not a transfer: it does
not move with any aggregate state and therefore contributes nothing to any
response reported below.

This is a deliberate limitation and it should be read as one. The model is a
heterogeneous-agent economy at the steady state and in its welfare accounting,
and a representative-agent economy in its aggregate transmission. The
approximation is defensible because the object the household side has to supply
to the mechanism of interest is an aggregate deposit-supply schedule --- the
sovereign shock reaches households through wages, dividends and taxes, none of
which is distribution-dependent in this environment --- but a full \citet{krusell1998income}
treatment, in which the cross-sectional distribution is itself carried as an
aggregate state, would let precautionary saving respond to sovereign risk and is
the natural next step.


% =============================================================================
% INSERT: as a new subsection after \subsection{Financial Intermediation},
%         before \subsection{Default Event and Sovereign Risk}.
% =============================================================================
\subsection{The union deposit market}
\label{sec:union-deposits}

Prices are flexible, so there is no nominal interest rate to be common across the
union. What a monetary union does bind, in real terms, is the market in
intermediary liabilities: a euro of deposit is a euro of deposit wherever it is
placed. We therefore replace the two national deposit-market clearing conditions
with \emph{one} union-wide clearing condition together with real interest parity.
This is the flexible-price image of a single policy rate plus national inflation
differentials, and it is the standard single-traded-bond margin of the
international real business cycle literature.

\paragraph{Clearing and the cross-border position.} Deposits are own-good claims
paying national rates. Converting $\F$'s book into $\D$-goods at $p_t$ and
weighting by mass, union clearing is
\begin{equation}
\underbrace{A_{\D,t}P^{c}_{\D,t}+\varkappa\,p_t\,A_{\F,t}P^{c}_{\F,t}}
           _{\text{household saving}}
\;=\;
\underbrace{\mathrm{dep}_{\D,t}+\varkappa\,p_t\,\mathrm{dep}_{\F,t}}
           _{\text{intermediary funding need}},
\qquad
\mathrm{dep}_{X,t}=\mathcal{A}_{X,t}-n_{X,t}.
\label{eq:mkt-dep}
\end{equation}
Because \eqref{eq:mkt-dep} is one condition where there used to be two, the
national deposit markets no longer have to balance separately. The gap
\begin{equation}
V^{\rm flow}_{\D,t}\;\equiv\;A_{\D,t}P^{c}_{\D,t}-\mathrm{dep}_{\D,t}
\label{eq:nfa-flow}
\end{equation}
is the net cross-border deposit position of country $\D$: the amount by which
$\D$'s households fund $\F$'s intermediaries rather than their own. It is the
absorption margin of the model. Without it each country is forced into national
saving--investment balance every period, and the sovereign shock cannot move
$\D$'s deposit rate without immediately moving $\D$'s investment in the opposite
direction --- the comovement problem that a closed-economy reading of this
mechanism runs into.

\paragraph{Parity and stationarity.} Deposit legs are own-good claims, so the
no-arbitrage condition between them is a \emph{real} parity condition,
\begin{equation}
1+r_{\D,t}
\;=\;\bigl(1+r_{\F,t}\bigr)\frac{\E_t\bigl[p_{t+1}\bigr]}{p_t}
\;-\;\kappa_{\rm nfa}\,\frac{V^{\rm flow}_{\D,t}}{\bar Y_\D}.
\label{eq:uip}
\end{equation}
The first term is the parity itself. Imposing instead the literal equality
$r_{\D,t}=r_{\F,t}$ with own-good legs is \emph{not} an equilibrium: the
real-exchange-rate valuation profit on the cross-border leg would be assigned to
nobody and the resource constraint would leak. The second term is the
debt-elastic premium of \citet{schmittgrohe2003closing}, the same device
\citet{bocola2016pass} uses to close his open economy. It is zero at the
symmetric steady state, so it is undistorting there; it is present because
without it the model has no force returning the external position to its
steady-state level, the cross-border position inherits a unit root, and an
impulse response that never returns cannot be read for a trough.
$\kappa_{\rm nfa}=0$ nests frictionless parity exactly.

\paragraph{The carried wealth state.} Households carry a gross deposit claim
$W_{X,t}$ and banks carry the gross obligation $P_{X,t}$ of \eqref{eq:P-state}.
Under national clearing these are the same object; under \eqref{eq:mkt-dep} they
differ, and \emph{one} state variable records the difference. Define
$V_t\equiv W_{\D,t}-P_{\D,t}$. Then
\begin{equation}
W_{\D,t}=P_{\D,t}+V_t,
\qquad
W_{\F,t}=P_{\F,t}-\frac{V_t}{\varkappa\,p_t},
\label{eq:union-wealth}
\end{equation}
so the union wealth identity
$W_{\D,t}+\varkappa p_t W_{\F,t}=P_{\D,t}+\varkappa p_t P_{\F,t}$ holds by
construction rather than as an imposed condition on separately approximated
states. The position accumulates at the $\D$ deposit rate,
\begin{equation}
V_{t+1}=\bigl(1+r_{\D,t}\bigr)\,V^{\rm flow}_{\D,t},
\qquad
W_{\D,t+1}=\bigl(1+r_{\D,t}\bigr)A_{\D,t}P^{c}_{\D,t}
           -\bigl(1+r^{wc}_{\D,t}\bigr)L_{\D,t},
\label{eq:V-lom}
\end{equation}
where the working-capital receivable is netted from the household's carried claim
for the same reason it is netted from the bank's obligation in
\eqref{eq:P-state}: the deduction is common to both legs of $V$, cancels
exactly, and leaving it out of one leg alone makes $W_\D$ and $P_\D$ drift apart
by $(1+r^{wc}_{\D,t})L_{\D,t}$ every period.


% =============================================================================
% INSERT: as a new subsection after \cref{sec:union-deposits}.
% =============================================================================
\subsection{Two sovereign markets, two bidders}
\label{sec:sovereign-markets}

Both sovereigns are held by both intermediaries, and in both markets the price
and the cross-border split are equilibrium objects rather than calibrated shares.
Under the single-$\lambda$ doctrine the banker has a first-order condition for
\emph{every} asset class, \eqref{eq:foc-general}, so each market has two demand
schedules. Specialising \eqref{eq:foc-general} to the $\D$ sovereign for each
bank, and writing $R_{X,t}=1+r_{X,t}$:
\begin{align}
\E_t\bigl[\Omega_{\D,t,t+1}\Xi^{\D}_{t+1}\bigr]
&=q^{\D}_{t}\Bigl(\E_t\bigl[\Omega_{\D,t,t+1}\bigr]R_{\D,t}+\lambda_\D\mu_{\D,t}\Bigr),
\label{eq:bondD-home}\\[3pt]
\E_t\bigl[\Omega_{\F,t,t+1}\Xi^{\D}_{t+1}\bigr]
&=q^{\D}_{t}\Bigl(\E_t\bigl[\Omega_{\F,t,t+1}\bigr]R_{\F,t}+\lambda_\F\mu_{\F,t}\Bigr)
  \Bigl[1+\psi_\F\frac{b^{\D}_{\F,t+1}-\bar b^{\D}_{\F}}{\bar B_\D}\Bigr],
\label{eq:bondD-foreign}
\end{align}
with the mirror-image pair pricing the $\F$ sovereign, in which the $\F$ bank
holds the home leg and the $\D$ bank the foreign one. Equation
\eqref{eq:bondD-home} determines the price $q^{\D}_t$, equation
\eqref{eq:bondD-foreign} the foreign holding $b^{\D}_{\F,t+1}$, and the residual
$b^{\D}_{\D,t+1}=B_{\D,t+1}-b^{\D}_{\F,t+1}$ clears the market by construction.

\paragraph{Why the foreign leg carries an adjustment cost.} The bracket in
\eqref{eq:bondD-foreign} is a portfolio adjustment cost on the cross-border
position, linear in the deviation from the steady-state holding and therefore
zero at the steady state: it does not distort the calibration. It is nonetheless
load-bearing. The two schedules \eqref{eq:bondD-home}--\eqref{eq:bondD-foreign}
differ only through $(\Omega_X,\mu_X,R_X)$, and a sovereign position is about
$7\%$ of an intermediary's divertible assets, so both demand curves are nearly
flat in the holding: the cross-border split is close to indeterminate and the
equilibrium is numerically ill-conditioned without $\psi$. The cost is what gives
the foreign schedule a definite slope. It is also the model's contagion channel:
it is because the $\D$ bank has a genuine first-order condition for the safe
$\F$ bond that a rise in $\pi^{d}$ produces a flight to quality that reprices
$q^{\F}$, rather than leaving the safe bond's price fixed by construction.

\paragraph{Two simplifications, stated.} The $\F$ sovereign's \emph{stock} is
held at its steady-state level, $B_{\F,t}\equiv\bar B_\F$: only its split between
the two intermediaries moves. The $\D$ sovereign's stock is fully endogenous,
evolving under \eqref{eq:debt-lom}--\eqref{eq:bohn-corrected} and absorbed by the
two banks each period. Clearing the $\D$ market against a fixed stock instead
opens a resource leak of about half a percent of GDP per five percent deviation
of the debt, which is why the debt is forward-integrated inside every equilibrium
evaluation rather than held fixed.


% =============================================================================
% INSERT: replacing / correcting eq. (bohn) in \subsection{Government}.
%         The surrounding text of that subsection is unchanged.
% =============================================================================
\subsection{The fiscal rule}
\label{sec:fiscal-rule-corrected}

Taxes respond to the \emph{surviving} debt stock $B_{X,t}h_t$ --- taxing the
pre-haircut stock in the default state would charge the household for a claim
that has just been written down --- and they do so linearly in the level:
\begin{equation}
T^{\tau}_{X,t}
=\bar T^{\tau}_X\bigl(\mathcal{B}_t\bigr)
 +\gamma_\tau\bigl(B_{X,t}h_t-\mathcal{B}_t\bigr),
\qquad
\bar T^{\tau}_X(\mathcal{B})=G_X+\delta_b\,\mathcal{B}\,\bigl(1-\bar q^{X}\bigr),
\label{eq:bohn-corrected}
\end{equation}
where the anchor $\mathcal{B}_t$ is the debt level the rule treats as normal:
$\mathcal{B}_t=\bar B_X$ in the no-default state and $\mathcal{B}_t=\varrho_D\bar
B_X$ in the default state. The coefficient $\gamma_\tau$ is not assigned. It is
inverted from the persistence of the debt stock the rule is required to deliver.
Substituting \eqref{eq:bohn-corrected} into \eqref{eq:debt-lom} and
differentiating with respect to the surviving stock $x=B_{X,t}h_t$,
\begin{equation}
\frac{\partial B_{X,t+1}}{\partial x}
=(1-\delta_b)+\frac{\delta_b-\gamma_\tau}{\bar q^{X}}\;\equiv\;\rho_B
\qquad\Longleftrightarrow\qquad
\gamma_\tau=\delta_b-\bigl[\rho_B-(1-\delta_b)\bigr]\bar q^{X},
\label{eq:gamma-tau}
\end{equation}
so that the calibrated object is the debt root $\rho_B$ --- a number with a
half-life one can read --- rather than a coefficient whose implied persistence
nobody inspects. At the anchor the rule returns $\bar T^{\tau}_X$ identically, so
the steady state is unchanged for any $\rho_B$ and no other block needs
recalibrating.

\paragraph{Why not the elasticity form.} Three readings of a unit fiscal
coefficient were tried, and \eqref{eq:bohn-corrected} is what survives. A unit
\emph{level} coefficient makes a $55\%$ haircut a fiscal windfall worth about
half of annual output and leaves the sovereign more indebted after default than
before it, which inverts the sign of the entire mechanism. A constant
\emph{elasticity}, $T^{\tau}=\bar T^{\tau}(B h/\bar B)^{\gamma_\tau}$ as in
\eqref{eq:bohn}, does not stabilise debt at all in this calibration: with
$G_X=0$ the steady-state tax is the net interest bill alone, about $0.3\%$ of
quarterly output against a debt stock of $98\%$ of it, so $\mathrm{d}T/\mathrm{d}B$
is three thousandths and the debt root is essentially unity --- the debt walks
out of any bounded approximation region within a few quarters of the shock.
Raising the elasticity to recover a root below one is worse: it swings the tax
level by a factor of ten across the ergodic range of the debt and by three orders
of magnitude at the default node, a convexity that a global polynomial basis
cannot carry, and the resulting approximation error lands on the steady state
itself. The rule linear in the level is exactly representable in the basis, and
its root is the object being calibrated.

\paragraph{The anchor must move with the regime.} If $\mathcal{B}_t$ were fixed at
$\bar B_X$ in both states, the term $\gamma_\tau(B_{X,t}h_t-\bar B_X)$ would
multiply the large negative deviation created by the write-down and hand the
household a tax cut worth several times the steady-state tax bill \emph{on
impact}, while the offsetting loss --- the haircut on the banks' books --- reaches
the household only slowly, through the exit-rate fraction $f$ of bank equity.
Default would then be a state the household is better off in, pricing more of it
would \emph{raise} consumption, and the risk premium in
\eqref{eq:bond-decomposition} would change sign. Re-anchoring the rule to the
post-default stock keeps the fiscal relief leg at roughly $0.2\%$ of output,
which is its correct magnitude. Because the default state is a discrete regime
with its own decision rules, the switch in $\mathcal{B}_t$ introduces no
kink into any equilibrium condition.


% =============================================================================
% INSERT: at the end of \subsection{Default Event and Sovereign Risk}, alongside
%         the s-process, so that both exogenous processes are stated together.
% =============================================================================
\paragraph{Productivity.} Total factor productivity in the periphery is the
second exogenous driver, and it is treated as a deterministic (perfect-foresight)
process: it is never innovated on the ergodic set, and enters only through the
transitional experiment that reads the economy's response along a decay path,
\begin{equation}
Z_{\D,t+1}=(1-\rho_z)\bar Z_\D+\rho_z Z_{\D,t},
\qquad Z_{\F,t}\equiv\bar Z_\F .
\label{eq:z-process}
\end{equation}
It is carried as a full state variable nonetheless, so that the
productivity experiment and the sovereign-risk experiment are read off the
\emph{same} solved decision rules and their responses are directly comparable.


% =============================================================================
% INSERT: as a new subsection after \subsection{Government}, i.e. the policy
%         block; it is the paper's application.
% =============================================================================
\subsection{A stochastic liquidity backstop}
\label{sec:ltro}

The policy experiment asks what a credible central-bank backstop does to an
economy in which sovereign risk is transmitted through intermediary balance
sheets. The instrument is the one the ECB actually deployed: collateralised
central-bank lending to banks, at the policy rate, against sovereign collateral.
With per-period probability $\pi^{\ell}$ the facility is offered; let
$\mathsf{cb}_t\in\{0,1\}$ record whether it is. The facility is a fixed envelope
of size $\ell_X$, and it is fully drawn whenever offered --- weakly optimal, since
it is lent at the rate the bank pays on deposits and relaxes a binding
constraint --- so there is no quantity to solve for and no complementarity
condition.

\paragraph{Why lending and not purchases.} It is worth establishing first that
the obvious alternative cannot work in this class of model.

\begin{proposition}[Liquidity ceiling]
\label{prop:ceiling}
Hold the continuation fixed. Any policy that operates on the sovereign bond price
\emph{through the intermediary's balance-sheet constraint alone} --- in
particular, any purchase of the outstanding stock --- cannot raise $q^{\D}_t$
above
\begin{equation}
\bar q^{\D}_t\;\equiv\;
\frac{\E_t\bigl[\Omega_{\D,t,t+1}\Xi^{\D}_{t+1}\bigr]}
     {\E_t\bigl[\Omega_{\D,t,t+1}\bigr]R_{\D,t}},
\label{eq:ceiling}
\end{equation}
the value of the same claim with the liquidity discount removed and nothing else.
\end{proposition}
\begin{proof}
Rearranging \eqref{eq:bondD-home}, $q^{\D}_t=\E_t[\Omega\Xi^{\D}]/(\E_t[\Omega]R_{\D,t}+\lambda_\D\mu_{\D,t})$,
which is strictly decreasing in $\mu_{\D,t}$; and $\mu_{\D,t}\ge0$ by
\eqref{eq:kkt}, with equality attainable. The supremum is therefore attained at
$\mu_{\D,t}=0$ and equals \eqref{eq:ceiling}.
\end{proof}

The ceiling is quantitatively tight, not merely logically binding: measured across
the state space of the solved model, the liquidity component of the $\D$ bond
price is between $0.2\%$ and $0.7\%$ of its level, and $0.63\%$ at the crisis
corner. A purchase programme that exhausted the entire liquidity discount would
therefore move the price by less than one percent, and in the solved model a
target of that kind absorbs a quarter to a third of the outstanding stock without
reaching it. Bond purchases in this economy do not work through the quantity they
remove. Central-bank \emph{lending} operates on a different margin and is not
subject to \cref{prop:ceiling}, because it moves the constraint on both sides at
once.

\paragraph{The mechanism is a single change to the incentive constraint.} Let the
facility be lent at the deposit rate, $r^{\ell}_{X,t}=r_{X,t}$, against the bank's
sovereign collateral. Central-bank credit is not divertible --- the collateral is
pledged --- so the constraint \eqref{eq:IC} becomes
\begin{equation}
\varphi_{X,t}\,\underbrace{\bigl(n_{X,t}+\mathsf{cb}_t\ell_X\bigr)}_{\textstyle n^{\rm IC}_{X,t}}
\;\ge\;
\lambda_X\underbrace{\bigl(\mathcal{A}_{X,t}-\mathsf{cb}_t\ell_X\bigr)}_{\textstyle \mathcal{A}^{\rm IC}_{X,t}},
\label{eq:IC-ltro}
\end{equation}
and \cref{prop:mu} goes through verbatim with $(n,\mathcal{A})$ replaced by
$(n^{\rm IC},\mathcal{A}^{\rm IC})$:
\begin{equation}
\mu_{X,t}=\max\left\{1-\frac{\E_t[\Omega_{X,t,t+1}]R_{X,t}\,n^{\rm IC}_{X,t}}
                            {\lambda_X\,\mathcal{A}^{\rm IC}_{X,t}},\;0\right\}.
\label{eq:mu-ltro}
\end{equation}
Portfolio shares, dividends and the net-worth recursion \eqref{eq:nw-lom}
continue to divide by \emph{actual} net worth $n_{X,t}$: $n^{\rm IC}$ is an object
in the diversion constraint, not equity.

Two margins move where a purchase moves one. The assets the facility funds leave
the divertible base, and the funding itself counts as equity in the constraint.
To a first order around any point,
\begin{equation}
\frac{\partial\bigl[n^{\rm IC}/\lambda\mathcal{A}^{\rm IC}\bigr]/\partial\ell
      \Big|_{\text{lending}}}
     {\partial\bigl[n/\lambda\mathcal{A}\bigr]/\partial\ell
      \Big|_{\text{purchase}}}
=1+\frac{\mathcal{A}_{X,t}}{n_{X,t}}
=1+\theta_{X,t},
\label{eq:relief-ratio}
\end{equation}
so at the calibrated leverage of five, a euro lent relieves the constraint six
times as much as a euro of bonds bought. The numerator term in
\eqref{eq:mu-ltro} is a margin no quantity of bond-buying can reach.

\begin{proposition}[Budget neutrality of the facility]
\label{prop:neutrality}
If $r^{\ell}_{X,t}=r_{X,t}$, then every flow budget identity of the model is
independent of $\mathsf{cb}_t$ and $\ell_X$. The entire effect of the facility
operates through \eqref{eq:IC-ltro}.
\end{proposition}
\begin{proof}
The balance sheet \eqref{eq:balance-sheet} becomes
$\mathcal{A}_{X,t}=n_{X,t}+\mathrm{dep}_{X,t}+\mathsf{cb}_t\ell_X$, so deposits are
$\mathrm{dep}_{X,t}=\mathcal{A}_{X,t}-n_{X,t}-\mathsf{cb}_t\ell_X$. The obligation
carried forward, \eqref{eq:P-state}, is then
\begin{align*}
P_{X,t+1}&=(1+r_{X,t})\mathrm{dep}_{X,t}
          +(1+r^{\ell}_{X,t})\mathsf{cb}_t\ell_X
          -(1+r^{wc}_{X,t})L_{X,t}\\
        &=(1+r_{X,t})\bigl(\mathcal{A}_{X,t}-n_{X,t}\bigr)
          -(1+r^{wc}_{X,t})L_{X,t},
\end{align*}
which is \eqref{eq:P-state} unchanged. On the household side a euro of deposit is
replaced by a euro of central-bank claim at the same rate, so
\eqref{eq:mkt-dep}, \eqref{eq:nfa-flow} and \eqref{eq:V-lom} are untouched. The
central bank lends at the rate it pays and bears no credit risk, so its carry is
identically zero and no remittance to the treasury is required.
\end{proof}

\cref{prop:neutrality} is what makes the experiment clean: because no resource
flow moves, any measured effect is attributable to the constraint, and the
redundant goods-market residual can be used as a test --- it must be unchanged to
machine precision between $\ell_X=0$ and $\ell_X>0$.

\paragraph{Three channels, one of which runs backwards.} The sign of the policy's
effect is not obvious ex ante, and the model is set up to decompose it rather than
to assert it. \emph{(a)} In the relieved state $\mu$ falls, so by \eqref{eq:rwc}
the credit spread falls and, through \eqref{eq:labour-demand}, hours and output
rise directly; and by \eqref{eq:bond-pricing} the liquidity discount on the
sovereign shrinks. \emph{(b)} The facility lowers $\mu'$, hence $\varphi'$, hence
$\Omega'$, \emph{most in exactly those states where the constraint is tightest} ---
which are the states in which $\Xi^{\D}$ is lowest. The covariance term in
\eqref{eq:bond-decomposition} shrinks, the risk premium falls, and the bond price
rises \emph{in every state, including those in which the facility is not offered}.
This is the announcement channel, and it is what allows a backstop that is never
drawn to stabilise the economy. It falls out of the quadrature; nothing is
imposed. \emph{(c)} Against these, the franchise-value channel runs the other way.
Since $\varphi=\E[\Omega]R/(1-\mu)$ and $\Omega=\beta[f+(1-f)\varphi']$, a safer
future lowers $\varphi'$, lowers $\E_t[\Omega]$, and by \eqref{eq:mu-closed}
\emph{raises} $\mu$ today: the bank's charter value is its collateral, so making
the future safer gives it less to lose. The channel is first-order here, because
$\mu$ is a small difference of numbers near one. Which of (a)+(b) and (c)
dominates is a general-equilibrium question, and a finding either way is a
result about standing liquidity backstops rather than a failure of the
experiment.

\paragraph{Availability in the default state, and size.} The facility is
available in the default state as well as outside it. This is the empirically
relevant configuration for an instrument that supports \emph{banks} rather than
the sovereign, and it is also what channel (b) requires: the default branch
carries little probability mass but by far the largest payoff deviation, so it
dominates $\mathrm{Cov}_t(\Omega,\Xi^{\D})$, the object a credible backstop
compresses. Withdrawing the facility in default --- historically accurate for
Greek collateral in 2012 and 2015 --- removes the largest single term of the
announcement channel and produces a collateral cliff; that is a separate
experiment, not a cheaper version of this one.

Size is a calibration decision with a trap in it, and it is resolved in favour of
a facility that \emph{relieves} the constraint without unbinding it. In this
calibration a facility of $2\%$ of quarterly output already drives $\mu$ to zero
at the steady state and $3.4\%$ drives it to zero in the crisis state; the
$40\%$ envelope used in \citet{bocola2016pass}'s own LTRO exercise is roughly
twelve times what is needed to neutralise the crisis entirely. At such a size the
whole relieved regime sits on the kink of $\max\{\cdot,0\}$, precisely the region
in which a global polynomial approximation of a $C^0$ multiplier is least reliable
(\cref{sec:kink}). The baseline sets $\ell_\D$ to halve the crisis-state
multiplier --- about $1.2\%$ of quarterly output --- which keeps $\mu>0$ in both
regimes; the large envelope is reported as a bounded variant with the caveat
attached. The facility is per country, with $\ell_\F=0$ for the targeted
experiment and $\ell_\F=\ell_\D$ for the union-wide one.

\paragraph{The compound regime.} $\pi^{\ell}$ is a policy scalar, not a state:
each solve is exact at its own $\pi^{\ell}$, and $\pi^{\ell}=0$ returns the
no-backstop economy exactly rather than approximately. Because the facility is a
second discrete event on which the continuation must be conditioned, the regime
index becomes the pair $(d_t,\mathsf{cb}_t)$, and \eqref{eq:expectation}
generalises to a product measure over the two independent regimes and the risk
innovation: for any function $g$ of next period's state,
\begin{equation}
\E_t[g]=\sum_{d'\in\{0,1\}}\;\sum_{\mathsf{cb}'\in\{0,1\}}
   \mathbb{P}_t(d')\,\mathbb{P}(\mathsf{cb}')\;
   \E^{s}_t\bigl[g(d',\mathsf{cb}',s_{t+1})\bigr],
\label{eq:expectation-compound}
\end{equation}
with $\mathbb{P}_t(d'=1)=\pi^{d}_t$ and
$\mathbb{P}(\mathsf{cb}'=1)=\pi^{\ell}$. Equation
\eqref{eq:expectation-compound} contains \eqref{eq:expectation} as the case
$\pi^{\ell}=0$.


% =============================================================================
% INSERT: closing subsection of \section{The Model}, replacing the sentence
%         "We relegate the description of the equilibrium conditions and market
%          clearing to the appendix." with a pointer to \cref{app:equilibrium}.
% =============================================================================
\subsection{Recursive competitive equilibrium}
\label{sec:rce}

The model is solved as a genuine recursive equilibrium, not as a perfect-foresight
transition, because the mechanism of interest is a risk premium: it is a property
of the conditional distribution of next period's states, and it is identically
zero in any certainty-equivalent approximation. The aggregate state is the
ten-dimensional vector
\begin{equation}
\mathcal{S}_t=\Bigl[
K_{\D,t},\;K_{\F,t},\;
P_{\D,t},\;P_{\F,t},\;
b^{\D}_{\D,t},\;b^{\D}_{\F,t},\;b^{\F}_{\D,t},\;
V_t,\;s_t,\;Z_{\D,t}\Bigr],
\label{eq:state}
\end{equation}
namely the two predetermined capital stocks \eqref{eq:capital-lom}, the two
intermediaries' carried obligations \eqref{eq:P-state}, the three carried
sovereign positions --- the fourth, $b^{\F}_{\F,t}$, is recovered from
\eqref{eq:mkt-bF} because $\bar B_\F$ is fixed, and $B_{\D,t}$ is recovered as
$b^{\D}_{\D,t}+b^{\D}_{\F,t}$ --- the cross-border deposit position
\eqref{eq:union-wealth}, the sovereign-risk factor \eqref{eq:s-process} and
productivity \eqref{eq:z-process}. Note that the central-bank facility adds
\emph{no} state, by \cref{prop:neutrality}: it carries no stock forward.

\begin{definition}[Recursive competitive equilibrium]
\label{def:rce}
A recursive competitive equilibrium is a set of decision rules, one for each
regime $j\leftrightarrow(d,\mathsf{cb})$,
\[
\Bigl\{N_{X},\,K'_{X},\,r_{X},\,A_{X},\,p,\,q^{\D},\,q^{\F},\,
b^{\D}_{\F}{}',\,b^{\F}_{\D}{}',\,\varphi_{X},\,C_{X},\,r^{wc}_{X}\Bigr\}(j,\mathcal{S}),
\]
together with the implied multipliers $\mu_X(j,\mathcal{S})$ from
\eqref{eq:mu-ltro} and the law of motion for $\mathcal{S}$ generated by
\eqref{eq:capital-lom}, \eqref{eq:P-state}, \eqref{eq:debt-lom},
\eqref{eq:V-lom}, \eqref{eq:s-process} and \eqref{eq:z-process}, such that at
every $(j,\mathcal{S})$: households solve
\eqref{eq:hh-objective}--\eqref{eq:hh-borrowing} in the aggregate form
\eqref{eq:agg-budget}--\eqref{eq:agg-euler}; firms and capital producers satisfy
\eqref{eq:labour-demand}--\eqref{eq:tobins-q}; intermediaries satisfy
\eqref{eq:foc-general} for every asset class together with the complementarity
\eqref{eq:kkt} in the closed form \eqref{eq:mu-ltro}; the government satisfies
\eqref{eq:govt-budget} and \eqref{eq:bohn-corrected}; expectations are taken
under \eqref{eq:expectation-compound}; and the markets of
\cref{app:equilibrium} clear.
\end{definition}

The equilibrium conditions, the market-clearing conditions and the redundancies
among them are collected in \cref{app:equilibrium}; the approximation of
\cref{def:rce} is described in \cref{app:solution}.


% =============================================================================
% INSERT: appendix, first section.
% =============================================================================
\section{Equilibrium conditions and market clearing}
\label{app:equilibrium}

\paragraph{Market clearing.} All quantities are per capita of their own country;
$\varkappa=\mathsf{M}_\F/\mathsf{M}_\D$ as in \cref{sec:size}.
\begin{align}
\text{$\D$ goods:}\quad
& Y_{\D,t}=P^{c}_{\D,t}C_{\D,t}+I_{\D,t}+G_\D+NX_{\D,t},
\label{eq:mkt-goods}\\
\text{$\F$ goods:}\quad
& Y_{\F,t}=P^{c}_{\F,t}C_{\F,t}+I_{\F,t}+G_\F+NX_{\F,t},
\label{eq:mkt-goods-F}\\
\text{$\D$ sovereign:}\quad
& b^{\D}_{\D,t+1}+b^{\D}_{\F,t+1}=B_{\D,t+1},
\label{eq:mkt-bD}\\
\text{$\F$ sovereign:}\quad
& b^{\F}_{\F,t+1}+\frac{b^{\F}_{\D,t+1}}{\varkappa}=\bar B_\F,
\label{eq:mkt-bF}\\
\text{union deposits:}\quad
& A_{\D,t}P^{c}_{\D,t}+\varkappa p_t A_{\F,t}P^{c}_{\F,t}
   =\mathrm{dep}_{\D,t}+\varkappa p_t\,\mathrm{dep}_{\F,t},
\label{eq:mkt-dep-repeat}\\
\text{capital and labour:}\quad
& K_{X,t+1}\ \text{held entirely by the country-$X$ intermediary},\quad
  N_{X,t}\ \text{clears \eqref{eq:labour-demand}.}
\label{eq:mkt-KN}
\end{align}

\paragraph{The solved system.} At each state and regime the equilibrium is the
root of thirteen conditions in the thirteen unknowns
\[
\bigl(N_\D,\;N_\F,\;K'_\D,\;K'_\F,\;r_\D,\;r_\F,\;p,\;
q^{\D},\;b^{\D}_{\F}{}',\;q^{\F},\;b^{\F}_{\D}{}',\;A_\D,\;A_\F\bigr).
\]
The system is solved jointly; the association below is indicative of which
condition principally determines which unknown, not of a recursive structure.

\begin{center}
\begin{tabular}{clc}
\hline
\# & Condition & principally determines\\
\hline
1--2   & capital Euler \eqref{eq:capital-euler}, $\D$ and $\F$ & $K'_\D,\;K'_\F$\\
3--4   & GHH intratemporal condition against \eqref{eq:labour-demand} & $N_\D,\;N_\F$\\
5      & $\D$ deposit Euler \eqref{eq:agg-euler} & $A_\D$\\
6      & deposit parity \eqref{eq:uip} & $r_\F$\\
7      & $\D$ goods market \eqref{eq:mkt-goods} & $p$\\
8      & $\D$ bank's $\D$-bond FOC \eqref{eq:bondD-home} & $q^{\D}$\\
9      & $\F$ bank's $\D$-bond FOC \eqref{eq:bondD-foreign} & $b^{\D}_{\F}{}'$\\
10     & $\F$ deposit Euler \eqref{eq:agg-euler} & $r_\D$\\
11     & union deposit clearing \eqref{eq:mkt-dep} & $A_\F$\\
12     & $\F$ bank's $\F$-bond FOC & $q^{\F}$\\
13     & $\D$ bank's $\F$-bond FOC & $b^{\F}_{\D}{}'$\\
\hline
\end{tabular}
\end{center}

Six further objects are determined by their own recursions rather than by market
clearing --- the two franchise values $\varphi_X$ from \eqref{eq:phi-recursion},
the two aggregate consumptions $C_X$ from \eqref{eq:agg-budget}, and the two
working-capital rates $r^{wc}_X$ from \eqref{eq:rwc}. In the global solution these
are carried as unknowns as well, each with the identity residual
$\log(\text{guess}/\text{implied})$, so that the system rooted at every
collocation node has nineteen equations per regime
(\cref{app:solution}).

\paragraph{Redundancy.} Given \eqref{eq:mkt-goods}, \eqref{eq:mkt-bD},
\eqref{eq:mkt-bF}, \eqref{eq:mkt-dep} and the budget constraints of households,
firms, intermediaries and the two governments, the $\F$ goods market
\eqref{eq:mkt-goods-F} and the current account are implied by Walras's law. They
are therefore dropped from the solved system and monitored as diagnostics: they
are the model's own internal accounting check, and reporting their size is how a
leak in any budget identity is detected. Under \cref{prop:neutrality} they
provide the sharpest available test of the policy experiment, since they must be
invariant to the facility to machine precision.

\paragraph{Timing conventions.} Three conventions are used throughout and each is
load-bearing. Capital producing at $t$ was purchased at $t-1$
\eqref{eq:production}, so a sovereign-risk shock cannot raise impact output
through investment and reaches output through hours alone. The deposit rate paid
at $t$ was set at $t-1$, in the bank's funding leg \eqref{eq:P-state}, in the
household's realised return \eqref{eq:hh-budget} and in the multiplier
\eqref{eq:mu-ltro}. And the working-capital loan is intra-period: it is extended
at $t$ at the rate \eqref{eq:rwc} set at $t$, it sits inside the divertible asset
base at the same $\lambda_X$ as every other class, and its repayment accrues to
bank net worth through \eqref{eq:P-state} --- not to households as a dividend,
which would make the credit spread a pure intra-period transfer with, under GHH
preferences, no wealth effect to offset it, and would turn the entire mechanism
expansionary.


% =============================================================================
% INSERT: appendix, second section.
% =============================================================================
\section{Solving the model}
\label{app:solution}

\paragraph{Approximation.} Each of the nineteen equilibrium objects of
\cref{app:equilibrium} is approximated, separately in each of the four regimes
$(d,\mathsf{cb})$, by a Chebyshev polynomial on a bounded box in the state
\eqref{eq:state}. The collocation grid is the Cartesian product of a Smolyak
sparse grid \citep{smolyak1963,kruegerkubler2004,judd2014smolyak} of level one
over the nine smooth dimensions and a dense Chebyshev factor of $m$ nodes in the
risk state $s$. The anisotropy is deliberate. All of the curvature in the problem
is in $s$, through the logistic \eqref{eq:pd}; the remaining dimensions are close
to linear over their ergodic range. Raising the Smolyak level buys resolution in
all ten dimensions at once, whereas the dense factor buys degree $m-1$ in the one
dimension that needs it, at full interaction with the sparse basis. Measured on a
test function with this model's curvature profile, the isotropic level-one grid
($21$ points) has relative root-mean-square error $1.9\times10^{-1}$ and the
isotropic level-two grid ($221$ points) $3.9\times10^{-2}$, against
$2.5\times10^{-2}$ for the refined grid at $m=5$ ($95$ points) and
$1.1\times10^{-3}$ at $m=9$ ($171$ points). Policies that are positive by
construction are collocated in logs, and interest rates as gross rates, so that
positivity is a property of the approximation rather than something enforced by
clipping.

\paragraph{Expectations.} The conditional expectations in
\eqref{eq:foc-general}, \eqref{eq:bond-pricing} and \eqref{eq:agg-euler} are
evaluated by \eqref{eq:expectation-compound}: a seven-node Gauss--Hermite rule
over the innovation to \eqref{eq:s-process}, crossed with the four
$(d',\mathsf{cb}')$ cells. In each cell the continuation is that regime's own
decision rules evaluated at the state that regime's law of motion actually
reaches, and the stochastic discount factor \eqref{eq:branch-sdf} is evaluated
branch by branch on that regime's own continuation. A regime carrying zero
probability is never evaluated, which is what makes $\pi^{d}\equiv0$ and
$\pi^{\ell}=0$ nest the smaller models exactly rather than to solver tolerance.

\paragraph{The solve.} The unknowns are the policy \emph{values} at the
collocation nodes --- nineteen objects $\times$ four regimes $\times$ the number
of nodes. One residual evaluation fits the Chebyshev coefficients to the current
guess, walks the equilibrium conditions of \cref{app:equilibrium} at every node
with \emph{that} interpolant as the continuation, and returns the stacked
residual; the whole vector is then handed to a single damped Newton step with a
finite-difference Jacobian. Because square collocation interpolates its own
nodes exactly, no equilibrium condition is evaluated against a frozen or lagged
continuation: the solved object is a fixed point of the equilibrium system
itself. The solution is reached along a homotopy --- a coarse grid with no
default risk, then the default regime introduced by walking the haircut down to
its calibrated value, then the joint system, then the same solution seeded onto
the refined grid --- and the acceptance criterion at every stage is that the sum
of squared residuals be no larger than a uniform $10^{-9}$ per equation. That
floor is the arithmetic limit of the equilibrium conditions themselves, in which
$O(1)$ expectations are differenced down to $O(10^{-4})$ excess returns, and it is
four orders of magnitude below any economic signal in the experiment.

\paragraph{The stochastic rest point.} The deterministic steady state of
\cref{sec:calibration} is \emph{not} the point at which the solved stochastic
model rests. Solved with $\pi^{d}\equiv0$ the model does sit on the deterministic
steady state indefinitely, which verifies both the steady state and the solver;
but once risk is priced, precaution moves the ergodic centre --- output and
consumption slightly below, investment and bank net worth slightly above, and the
incentive constraint measurably slacker. \citet{bocola2016pass} reports the same
gap in his own solution and handles it the same way. Every impulse response in
this paper is therefore constructed as he constructs his: the economy is first
simulated without shocks until it settles at its own rest point $\mathcal{S}^{*}$;
the response is then the difference between a shocked path and an unshocked path
both started at $\mathcal{S}^{*}$. Reading a response against the deterministic
steady state instead charges the walk between the two rest points to the shock,
which in this model accounts for more than half of the apparent post-impact
dynamics.

\paragraph{The accuracy limit, and how it is reported.}
\label{sec:kink}
The multiplier \eqref{eq:mu-ltro} is continuous but not differentiable at
$\mu=0$, and the economy rests near that kink. A polynomial interpolant of a
$C^{0}$ function returns $\mu>0$ in a neighbourhood where the truth is zero, so
two legitimate reads of the same solved model --- evaluating the fitted rules at a
state, and clearing the equilibrium conditions exactly at that state against the
same continuation --- can differ by more than the response being measured. This is
a Gibbs phenomenon and it shrinks slowly with resolution. It is not peculiar to
this model: the same measurement on \citet{bocola2016pass}'s own published
solution gives a liquidity premium of $28$ basis points from his fitted policy
against $2$ basis points from the exact multiplier, and the fitted number is the
one he reports. Two things are done about it here. The intermediation wedge is
calibrated so that the incentive constraint binds at the ergodic rest point
rather than resting exactly on the kink, which collapses the discrepancy by a
factor of four; and every reported number is accompanied by both reads, so that
the bracket is visible rather than implicit. Where the two reads disagree, the
pair is the honest object, and the \emph{level} of the response should be read as
bracketed rather than point-identified.
